\section{Model}\label{sec:model}

A stylized model captures what the set-aside does and, as important, what it does not do. $N$ potential bidders compete for a single indivisible contract. Each bidder $i$ has type $k_i \in \{\text{SME}, \neg\}$ and draws a private cost of supply $c_i$ from a type-specific distribution $F_c^k$, with $F_c^{\text{SME}}$ and $F_c^{\neg}$ differing in ways the data are allowed to determine.

The auction format is the identifying payload. BEC implements procurement through a Preg\~ao---a descending-clock electronic English-reverse auction. The clock opens at the buyer's reference price $p^{\mathrm{ref}}$, firms post downward-revisable offers until only one firm remains, and that firm supplies the contract. Under \citet{haile2003}'s independent-private-values result, the weakly dominant strategy is to exit at one's own cost, so drop-out prices of losers \emph{point-identify} their costs without inverting a bid function or imposing parametric structure. The winner leaves only an upper bound $c_{\text{win}} \le c_{(2)}$ recovered by the Turnbull NPMLE robustness in Section~\ref{sec:robustness}. The first-price-sealed-bid settings of \citet{marion2007}, \citet{krasnokutskaya2011}, and \citet{atheyseira2011} require cross-equation constraints to disentangle channels that the Preg\~ao format identifies directly.

The choice of asymmetric IPV with the Krasnokutskaya UH correction and Athey--Seira reduced-form entry is the lightest specification that sustains the within-auction-vs-entry decomposition. Three alternatives are bounded explicitly in the robustness battery: common-value or affiliated-private-value formulations \citep{milgrom1982, hong2003} are accommodated up to within-auction cost correlation $\rho_c \le \ipvRhoMax$ without rejecting the IPV restriction (Section~\ref{sec:robustness}); nested entry-bidding equilibria \`a la \citet{krasnokutskayaseim2011} are eschewed deliberately to keep the welfare statement anchored to the legal trigger rather than to a participation-cost identity it does not directly identify; open-format and FPSB equivalents are bounded under \citet{maskinriley2000} with the V3-vs-V0 ranking preserved at the upper bound (Section~\ref{sec:fpsb_bound}).

\begin{assumption}[Independent Private Values, asymmetric across types]\label{a:ipv}
Costs $c_i$ are drawn independently across bidders. Conditional on type $k$, $c_i \sim F_c^k$ with density $f_c^k$ on support $[\underline{c}, \overline{c}] \subset \mathbb{R}_+$. The SME and non-SME distributions are not constrained to coincide.
\end{assumption}

The empirical median-cost gap of \medCostGapLo--\medCostGapHi{} of the reference price across (pharma $\times$ type) cells supports asymmetry directly; symmetric IPV would require this gap to vanish.

\begin{assumption}[Multiplicative auction-level heterogeneity]\label{a:uh}
The normalized log-bid decomposes as $y_{it} = \log(b_{it} / p_t^{\mathrm{ref}}) = a_t + e_{it}$, where $a_t$ is an auction-level scale common to all bidders in auction $t$ and $e_{it}$ is bidder-specific and independent across bidders within $t$. Both are mean-zero with finite variances $\sigma_a^{2,k}$ and $\sigma_e^{2,k}$ that may vary by stratum.
\end{assumption}

Assumption~\ref{a:uh} is the \citet{krasnokutskaya2011} restriction; multiplicativity is the natural form when item complexity or ANVISA burden scales costs uniformly across firms. The Turnbull NPMLE corroborates the baseline within \turnbullDeviationPct~percent of the total effect, dispensing with Assumption~\ref{a:uh} entirely.

\begin{assumption}[Set-aside, with equilibrium SME selection]\label{a:setaside}
Under the open regime both SMEs and non-SMEs enter with strictly positive probability. Under the SME-only regime, only SMEs are admissible. Within a short window around the cutoff the underlying technology of supply is stable, but the equilibrium distribution of active SMEs under the restricted regime may differ from its pre-regime counterpart because the policy changes which SMEs find entry profitable. Write these distributions $F_c^{\text{SME},\tau}$ for $\tau \in \{\text{Pre},\text{Post}\}$.
\end{assumption}

\begin{proposition}[Identification of $F_c^k$ under ascending-clock IPV]\label{prop:identification}
Under Assumptions~\ref{a:ipv}--\ref{a:setaside}, drop-out exit prices point-identify the type-specific cost distribution $F_c^k$ \citep{haile2003}, up to the auction-level scale $a_t$ recovered by the variance decomposition in Assumption~\ref{a:uh} \citep{krasnokutskaya2011}. The post-policy SME distribution $F_c^{\text{SME,Post}}$ is identified as an equilibrium-selection object; strict invariance ($F_c^{\text{SME,Post}} = F_c^{\text{SME,Pre}}$) is recovered as a nested benchmark.
\end{proposition}

Assumption~\ref{a:setaside} is the structural counterpart of parallel trends, with equilibrium selection made explicit \`a la \citet{atheyseira2011}: the policy changes the admissible set and therefore the composition of SME bidders who enter, so $F_c^{\text{SME,Post}}$ is an equilibrium object. The primitive-invariance KS test (Table~\ref{tab:v3_primitive_invariance}) rejects strict invariance of non-SME $F_c$ in three of four Preg\~ao strata after UH cleaning (KS distances \ksNpPregaoUH{} non-pharma, \ksPharmaPregaoUH{} pharma), so a structural exercise that imposed strict invariance on data that visibly reject it would be misspecified by construction. The main specification therefore adopts equilibrium selection and reports strict invariance as a limiting bound in Section~\ref{sec:robustness}.

\paragraph{Counterfactual simulation.} Under ascending-clock IPV, exit-at-cost is weakly dominant, the surviving bidder pays the second-lowest exit, and the expected transaction price equals $\mathbb{E}[c_{(2)}]$; revenue equivalence \citep{vickrey1961, maskinriley2000} extends the same object to FPSB. The Bayes--Nash counterfactual reduces to Monte Carlo simulation over draws from the estimated $F_c^k$ at the type-specific Poisson arrival rates, with observed Pre counts $(N^{\text{SME}}_{\text{Pre}}, N^{\neg}_{\text{Pre}})$ as status-quo arrivals and Post counts $(N^{\text{SME}}_{\text{Post}}, 0)$ as equilibrium arrival under the restricted regime in the \citet{atheyseira2011} spirit. The implied zero-profit entry costs (Table~\ref{tab:v3_entry_cost}, \entryCostMin--\entryCostMax{} per bidder) are per-auction calibration margins, not jointly identified fixed costs of the kind in \citet{krasnokutskayaseim2011}.

\paragraph{The decomposition.} Three counterfactual scenarios anchor the exercise: $S_1$ is the open regime at the pre-period pool (baseline); $S_2$ imposes SME-only at the pre-period SME pool size, isolating the within-auction component; $S_3$ uses the observed post-policy arrival rate and SME cost distribution, the full policy effect with endogenous entry. The simulated total effect decomposes arithmetically as
\begin{equation}
\underbrace{p_{S_3} - p_{S_1}}_{\text{simulated total effect}}
\;=\; \underbrace{(p_{S_2} - p_{S_1})}_{\text{within-auction (sheltered bidding)}}
   \,+\, \underbrace{(p_{S_3} - p_{S_2})}_{\text{participation margin}}.
\end{equation}
The within-auction component can exceed the realized total effect when participation partially offsets it---which the data deliver in non-pharmaceutical markets.

\paragraph{Sheltered bidding as a characterized property.} The within-auction component generalizes beyond the BEC-SP application: any restriction of admissibility to a single bidder type, holding the on-type pool fixed, generates an order-statistic shift identified from $F_c^k$ alone.

\begin{definition}[Sheltered bidding]\label{def:sheltered}
Consider an asymmetric IPV auction in which the buyer restricts admissibility to bidder type $k$. Let $p_{\text{open}}$ and $p_{\text{shelter}}$ denote the expected transaction prices under the open and type-$k$-only regimes at the same on-type pool size. The \emph{sheltered-bidding component} is $\Delta_{\text{shelter}} \equiv p_{\text{shelter}} - p_{\text{open}}$, identified separately from the \emph{participation-margin component} $\Delta_{\text{entry}} \equiv p_{\text{endo}} - p_{\text{shelter}}$ that captures the equilibrium response of type-$k$ entry. In BEC-SP notation, $\Delta_{\text{shelter}} = p_{S_2} - p_{S_1}$ and $\Delta_{\text{entry}} = p_{S_3} - p_{S_2}$.
\end{definition}

Under Assumptions~\ref{a:ipv}--\ref{a:setaside}, $\Delta_{\text{shelter}}$ is point-identified from drop-out exits via Proposition~\ref{prop:identification}; the expected second-order statistic shifts mechanically as a function of pool composition, without any active behavioral adjustment by surviving bidders. The component is structural in the sense that any policy intervention removing a bidder type while holding the surviving pool fixed produces a $\Delta_{\text{shelter}}$ predictable from $F_c^k$ and $F_c^{\neg}$ alone: procurement set-asides, exclusionary clearance rules, supplier disqualification regimes. The decomposition predicts $|\Delta_{\text{shelter}}| > |\Delta_{\text{entry}}|$ when the off-type pool is thick enough to materially shift the second-order statistic and the on-type entry response is bounded. Both conditions are satisfied in the thick standardized non-pharmaceutical case; the pharmaceutical case sits closer to the boundary on both margins, and the bifurcation between the model-stable non-pharma decomposition and the model-sensitive pharma one is a prediction of the framework rather than an artifact of the application.
